\section{Modern KG reasoning: multi-hop reasoning (Option C)}\label{sec:multihop}

\subsection{Background and research}
Multi-hop reasoning answers a question by following several connected relationships, so the answer is the end
of a \emph{path} and the path itself is the explanation. Lao et al.\ \citep{lao2011pra} showed that sequences of
relations (relation paths) are informative features for inference in large knowledge bases; MINERVA
\citep{das2018minerva} learned to walk from a query entity to the answer, returning the walk as evidence; and
Think-on-Graph \citep{sun2024tog} lets an LLM explore graph paths step by step so that answers stay
traceable. Hogan et al.\ \citep{hogan2021knowledge} place these methods alongside deductive (rule-based)
reasoning and stress that graph paths make conclusions explainable.

\textbf{What we took from the literature.} From path-ranking work we took the idea that each question type is
a \emph{relation-path pattern} (for example \texttt{participatesIn} $\rightarrow$ \texttt{hasParticipant}
$\rightarrow$ \texttt{standingOf}$^{-1}$). From MINERVA and Think-on-Graph we took the principle that the
returned path, with every intermediate entity, is part of the answer. We did not train a learned walker:
our graph is small (about four thousand triples) and the answers must be exactly right, so deterministic graph
traversal over relation patterns is both sufficient and verifiable, while learned walkers are designed for
large, incomplete graphs.

\subsection{Method}
The reasoner (\texttt{qualifypath/multihop.py}) works on the graph after OWL~2~RL reasoning and rules R1--R7.
\begin{enumerate}
  \item \textbf{Generic path search} --- breadth-first search over entity-to-entity relations
        (\texttt{participatesIn}, \texttt{hasWinner}, \texttt{hasLoser}, \texttt{playedAt},
        \texttt{scheduledAt}, \texttt{locatedIn}, \texttt{homeVenue}), walking each relation forwards or
        backwards and returning all shortest connecting paths (e.g.\ RR to Wankhede Stadium).
  \item \textbf{Question-specific chains} for CQ4 (remaining fixtures), CQ5--CQ7 (qualification analysis,
        including wins needed to stay alive and to guarantee a place), CQ8 (direct rivals), CQ10 (what-if),
        CQ11 (teams beaten and where), CQ13 (venues via the inferred chain) and transitive win chains
        (A beat B, B beat C).
  \item \textbf{Provenance on every hop} --- each hop is checked against the three graph layers and labelled
        \emph{explicit} (stored), \emph{inferred} (OWL~2~RL) or \emph{rule-derived} (R1--R7).
\end{enumerate}

\textbf{Predicted vs.\ verified facts.} Option~C makes no statistical predictions: every hop is a stored or
logically inferred fact. The only hypothetical output is the CQ10 what-if analysis, which is always labelled as
a \emph{scenario} and never presented as a fact.

\textbf{Why it suits the domain.} Qualification questions are inherently multi-hop: a team's chances depend on
its remaining matches, their opponents and those opponents' standings. Showing the full path is exactly the
explanation a fan or analyst needs.

\subsection{Representative inputs and outputs}
\lstinputlisting[basicstyle=\ttfamily\tiny,caption={Reasoner output: question, answer, reasoning steps and
labelled paths (excerpt of \texttt{Results/iteration2/multihop\_examples.txt}).}]{generated/multihop_examples.txt}

\shot[0.9\linewidth]{app_reasoner_lsg.jpg}{Multi-hop reasoner page: qualification analysis for LSG, with
reasoning steps, labelled paths and the path graph.}
\shot[0.9\linewidth]{app_reasoner_rr_venues.jpg}{Venues RR has played at: each answer is the inferred
\texttt{playedAtVenue} link plus the match that supports it.}

\subsection{Evaluation}
Ground truth was recomputed with plain Python from the cleaned match table, without RDF, SPARQL or the
reasoner, so the two implementations share no code. Every team was queried for each reasoning type and every
ordered pair of teams for win chains (150 questions). We also checked that every hop of every returned path
exists in the graph.

\begin{table}[H]\centering
{\small
\begin{tabular}{lrrrr}
\toprule
\textbf{Reasoning type} & \textbf{Questions} & \textbf{Correct} & \textbf{Accuracy} & \textbf{Paths} \\
\midrule
CQ4 remaining fixtures & 10 & 10 & 100\% & 62 \\
CQ11 teams beaten & 10 & 10 & 100\% & 38 \\
CQ13 venues played (chain) & 10 & 10 & 100\% & 50 \\
CQ5-7 qualification & 10 & 10 & 100\% & 10 \\
CQ8 rivals & 10 & 10 & 100\% & 62 \\
CQ10 what-if (scenario) & 10 & 10 & 100\% & 10 \\
Win chains (\textless{}=3 wins) & 90 & 90 & 100\% & 457 \\
Total & 150 & 150 & 100.0\% & 689 \\
\bottomrule
\end{tabular}}

\caption{Multi-hop reasoning accuracy against independent ground truth.}
\end{table}
\begin{itemize}\setlength\itemsep{1pt}
\item Hops checked: 3356; hops not found in the graph: 0.
\item Average path length: 4.81 hops (max 6).
\item Hop provenance: explicit: 2948 (88\%), inferred (OWL 2 RL): 244 (7\%), rule-derived (SPARQL R1-R7): 164 (5\%)
\item CQ10 answers are hypothetical scenarios and are labelled as such, never as facts.
\end{itemize}


All 150 answers match, and no path contains a hop that is absent from the graph. In a first run, 27 win-chain
answers listed the same team chain twice because SRH beat RR in two different matches; we now report unique team
chains while keeping both supporting match paths. Because the ground truth uses the same cleaned data, this
evaluation verifies the reasoning, not the data; the data itself was validated against the official points table
(Section~2).
